% CONDITIONAL SECTION — included under Option A of paper/PLAN.md §1.3 (scope
% ruling still open for T0). Under Option B this file is deleted wholesale,
% together with 08-dolgachev-doran.tex; no other section depends on it except
% via the labels sec:lattice, nc:monodromy, prop:usplit.
% ============================================================================
% 06-lattice.tex — monodromy and the integral lattice. Claims #14-#17 of
% paper/PLAN.md §3. Every number traces to C2_cooper_s7_v4.json (SHA256
% 036cd895db892aee9802514ec18668535f8896cee53cf6123533af9844387c3e).
% ============================================================================
\section{The monodromy lattice of the \texorpdfstring{$s_7$}{s7} family}
\label{sec:lattice}

\subsection*{Disclosure}

\emph{The monodromy matrices used in this section are obtained numerically.}
They are computed by analytic continuation at 60-digit working precision and
then recognized as exact algebraic numbers. That recognition step is not a
proof, and nothing derived from it is stated as a theorem here: the recognized
matrices appear as Numerical Claim~\ref{nc:monodromy}, status \textbf{(N)}.
Everything downstream of them --- the invariant form, the orbit lattice, the
Gram matrix, and the integral splitting --- is exact integer and rational
arithmetic, and is stated as a proposition \emph{about the recognized
matrices}, i.e. conditionally on Numerical Claim~\ref{nc:monodromy}. The one
statement in this section that is unconditional, because it is a finite
verifiable assertion about an explicitly displayed integer matrix, is
Proposition~\ref{prop:usplit}. We are equally explicit in
\S\ref{sec:framework} about what is \emph{not} claimed: nothing here
identifies the computed lattice with any invariant of a surface.

\subsection{Numerical continuation}

Work in the Frobenius basis of $\Ltwo$ at the MUM point $z=0$: the
holomorphic solution $A(z)$ with $A(0)=1$, and
\[
  \widehat{B}(z) \;=\; \frac{A(z)\log z + g(z)}{2\pi i},
  \qquad g \text{ holomorphic},\ g(0)=0,
\]
whose existence is Lemma~\ref{lem:unipotent}. In this basis the local
monodromy at $z=0$ is $\left[\begin{smallmatrix}1&1\\0&1\end{smallmatrix}\right]$
by construction.

The finite singular loci are computed at runtime as the roots of the leading
$\thetaop$-coefficient of $\Ltwo$, giving $\{-1,\tfrac1{27}\}$ for $s_7$; they
are not typed in. Continuation is by Taylor stepping at 60 decimal digits of
working precision, series order $140$, step ratio $0.35$, with the series tail
bounded by $10^{-45}$ at every step and the bound asserted rather than
assumed. Loops are taken based at $z_0 = \rho/4$, where $\rho$ is the distance
from $0$ to the nearest singular point.

\begin{remark}[Machinery control]\label{rem:machinery-control}
Before any unknown loop is computed, the numerical machinery is required to
reproduce a matrix known analytically: the continued loop around $z=0$ must
return $\left[\begin{smallmatrix}1&1\\0&1\end{smallmatrix}\right]$. It does, to
$3.17\times 10^{-61}$. A pipeline that cannot recover a known answer is not
used to produce unknown ones.
\end{remark}

\begin{numclaim}[Elliptic monodromies of the $s_7$ partner; status (N)]
\label{nc:monodromy}
In the Frobenius basis above, the monodromy matrices of $\Ltwo$ around the two
finite singular points are
\[
  M(\tfrac1{27}) \;=\; \frac{i}{\sqrt7}
  \begin{pmatrix} 0 & 1 \\ -7 & 0\end{pmatrix},
  \qquad
  M(-1) \;=\; \frac{i}{\sqrt7}
  \begin{pmatrix} 7 & 4 \\ -14 & -7\end{pmatrix} .
\]
\end{numclaim}

The evidence, and the limits of that evidence, are as follows. The two
numerically continued matrices satisfy $\det M = -1$ and $M^2 = \mathrm{Id}$ to
about $10^{-59}$, as the local exponents $\{0,\tfrac12\}$ of
Proposition~\ref{prop:L2-scheme} require. The quantity $\sqrt7$ is not put in
by hand: it emerges from the continuation. What is actually recognized is not
the $2\times 2$ matrices themselves but their symmetric squares in the flag
basis $(A^2, A\widehat{B}, \widehat{B}^2)$, which are rational; recognition is
performed with a loud acceptance gate of $10^{-35}$, the largest observed
residual is ${\approx}10^{-59}$, and the denominators appearing in the
recognized entries come out as $\{1,7\}$. The recognized matrices are then
subjected to exact post-verification: each is an involution; the product of
the three loops matches the monodromy at $\infty$ predicted by the
independently computed exponents $\{\tfrac13,\tfrac23\}$ there (trace $0$,
order $3$); the joint invariant symmetric bilinear form exists and is unique
up to scale; and the orbit lattice below closes under all generators and their
inverses. A conspiracy of numerical errors surviving all of these exact gates
is not credible. It is nevertheless not a proof, which is why this is a
Numerical Claim.

\subsection{The exact lattice pipeline}

From here on, all arithmetic is exact --- integers and rationals only.

\begin{proposition}[The joint monodromy-invariant lattice; status (E), given
Numerical Claim~\ref{nc:monodromy}]\label{prop:lattice}
Let $N(p) = \Symtwo M(p)$ denote the recognized rational $3\times3$ matrices
and let $\Gamma$ be the group they generate together with the cusp matrix
$N(0) = \left[\begin{smallmatrix}1&1&1\\0&1&2\\0&0&1\end{smallmatrix}\right]$.
Then:
\begin{enumerate}
\item the space of $\Gamma$-invariant symmetric bilinear forms on $\Q^3$ is
  one-dimensional;
\item the $\Gamma$-orbit of the cusp-invariant isotropic vector $f$ spans a
  lattice that stabilizes after finitely many steps and is closed under all
  generators and inverses;
\item in its primitive even rescaling, the restriction of the invariant form
  to that lattice has Gram matrix
  \[
    G \;=\;
    \begin{pmatrix} 0&0&-1\\ 0&14&0\\ -1&0&0 \end{pmatrix},
  \]
  with $\det G = -14$, signature $(2,1)$, discriminant group $\Z/14$ and
  discriminant form (in the sense of~\cite{Nikulin1979}) taking the value
  $q = \tfrac{1}{14} \bmod 2\Z$ on a generator;
\item the divisibility of $(N(0)-\mathrm{Id})^2$ on the lattice yields
  $2n = 14$;
\item the lattice admits no proper even $\Gamma$-invariant overlattice.
\end{enumerate}
\end{proposition}

Items (i)--(v) are finite exact computations on the recognized matrices; they
are performed by \artifact{check\_U1\_lattice.py}, whose structural conditions
are asserted (and therefore can fail) rather than reported. The sign of the
form is fixed by an independent positivity re-derivation of the framework
period shape, not chosen.

\subsection{The integral splitting}

The next statement is the one place in this section where nothing is
conditional. It is a finite assertion about the explicit integer matrix $G$
displayed above, checkable by hand in a minute, and we give the witness so
that the reader can do so.

\begin{proposition}[$G \cong \U\oplus\lat{14}$ over $\Z$; status (E),
unconditional]\label{prop:usplit}
Let $G$ be the integer matrix of Proposition~\ref{prop:lattice}(iii) and let
\[
  P \;=\;
  \begin{pmatrix} 1&0&0\\ 0&0&-1\\ 0&-1&0\end{pmatrix},
  \qquad \det P = -1 .
\]
Then $P \in \GL_3(\Z)$ and
\[
  P^{\mathsf T} G P \;=\;
  \begin{pmatrix} 0&1&0\\ 1&0&0\\ 0&0&14 \end{pmatrix}
  \;=\; \U \oplus \lat{14}.
\]
In particular the lattice with Gram matrix $G$ is isometric over $\Z$ to
$\U\oplus\lat{14}$.
\end{proposition}

\begin{proof}
Write $p_1,p_2,p_3$ for the columns of $P$ and $\langle\cdot,\cdot\rangle$ for
the form with Gram matrix $G$ in the standard basis $e_1,e_2,e_3$. Then
$p_1 = e_1$, $p_2 = -e_3$, $p_3 = -e_2$, and directly from $G$:
$\langle p_1,p_1\rangle = G_{11} = 0$;
$\langle p_1,p_2\rangle = -G_{13} = 1$;
$\langle p_1,p_3\rangle = -G_{12} = 0$;
$\langle p_2,p_2\rangle = G_{33} = 0$;
$\langle p_2,p_3\rangle = G_{32} = 0$;
$\langle p_3,p_3\rangle = G_{22} = 14$.
So $\{p_1,p_2\}$ spans a hyperbolic plane $\U$, orthogonal to $\langle
p_3\rangle \cong \lat{14}$. Expanding $\det P$ along the first column gives
$\det\left[\begin{smallmatrix}0&-1\\-1&0\end{smallmatrix}\right] = -1$, so $P$
is unimodular and the base change is an isometry of lattices over $\Z$.
\end{proof}

\begin{remark}[Two independent verifications]\label{rem:two-verifications}
Proposition~\ref{prop:usplit} has been verified by two implementations written
independently of one another and sharing no code.
\begin{enumerate}
\item The pipeline of \artifact{check\_U1\_lattice.py} constructs a base change
  of determinant $+1$ during its stage 3 and checks $S^{\mathsf T}GS =
  \U\oplus\lat{14}$ exactly. It serializes $\det S$ and $S^{\mathsf T}GS$ to
  its certificate but not $S$ itself.
\item A from-scratch implementation, \artifact{check\_U1\_splitting\_independent.py},
  loads the certificate at runtime, reads only $G$ and the integer $d$ (never
  hardcoding $14$ outside a labeled control), assembles the target
  $\U\oplus\lat{d}$ at runtime, and re-derives its own witness from $G$ alone
  --- locating a primitive isotropic vector of unimodular pairing, completing
  it to a hyperbolic pair by an extended-gcd combination, and taking the
  orthogonal-complement generator. All arithmetic is exact Python integers. It
  returns the matrix $P$ displayed in Proposition~\ref{prop:usplit}, of
  determinant $-1$.
\end{enumerate}
The two witnesses differ (determinants $+1$ and $-1$); that they differ is
expected, since any element of $\GL_3(\Z)$ realizing the isometry is an equally
valid proof, and it confirms that the second implementation did not reproduce
the first. The second is shipped with five negative controls: perturbing one
entry of $G$ must break agreement with the certificate's own determinant, and
does; corrupting the witness must break $P^{\mathsf T}GP = \U\oplus\lat{d}$,
and does; doubling a column of the witness must be rejected at the
$\GL_3(\Z)$ gate, and is; testing the genuine $s_7$ Gram against the target
$\U\oplus\lat{20}$ must fail, and does; and the unmodified certificate must
still pass after all scrambling hooks have been exercised, and does.
\end{remark}

\subsection{The \texorpdfstring{$s_{10}$}{s10} control}

\begin{numclaim}[Level discrimination; status (N)+(E)]\label{nc:s10-control}
The identical pipeline --- same code, same tolerances, same exact
post-verification --- applied to the $s_{10}$ family derives determinant $-20$,
$2n = 20$, and an integral splitting $\U\oplus\lat{20}$.
\end{numclaim}

This is the control that matters most. A pipeline that produced
$\U\oplus\lat{14}$ from any input, or that had $14$ wired into it anywhere,
would be worthless; running it unchanged on a different member of the same
template and obtaining a different, internally consistent answer at a
different level shows that the computation is sensitive to the operator it is
given. The certificate records the $s_{10}$ values as computed-versus-computed
inequalities against the $s_7$ values, with no expected result typed in.
Further controls in the same suite corrupt the recognized matrices in three
distinct ways: a raw entry perturbation fails at the involution gate; a
conjugated involution fails at the infinity-product gate; and a corruption fed
directly to the invariant-form step causes the joint invariant form to vanish
(solution space of dimension $0$) rather than to be silently replaced by a
different form.

\begin{remark}[What has and has not been established]
\label{rem:lattice-scope}
Established, conditionally on Numerical Claim~\ref{nc:monodromy}: the joint
monodromy-invariant lattice of the $s_7$ family, in its primitive even
scaling, is isometric over $\Z$ to $\U\oplus\lat{14}$, with the discriminant
data of Proposition~\ref{prop:lattice}. Established unconditionally:
Proposition~\ref{prop:usplit}, a statement about the displayed matrix $G$.
\emph{Not} established anywhere in this section: that this lattice is the
transcendental lattice of anything. That identification is a separate question,
taken up --- and left conditional --- in \S\ref{sec:framework}.
\end{remark}

% ----------------------------------------------------------------------------
% Generated-by: Opus 5 (Stream 1 paper agent) | Verified-by: all numbers read
% from Stream 2 data/certificates/C2_cooper_s7_v4.json (derived block, controls
% block, how block) and briefs/STREAM2_U1_EXECUTION_2026_07_27.md; witness P and
% control inventory from this repo's briefs/STREAM1_U1_INDEPENDENT_VERIFICATION_
% 2026_07_27.md and checkers/check_U1_splitting_independent.py; P^T G P verified
% by hand in the proof. Claims #14-#17 of paper/PLAN.md §3.
% | Reviewed-by: pending T0 (Xavier)
% ----------------------------------------------------------------------------
